\documentclass[11pt,a4paper]{article}
\usepackage{fontspec,amsmath,amssymb,amsthm,microtype}
\setmainfont{lmroman10-regular.otf}[BoldFont=lmroman10-bold.otf,ItalicFont=lmroman10-italic.otf,BoldItalicFont=lmroman10-bolditalic.otf]
\setmonofont{lmmono10-regular.otf}[Scale=MatchLowercase]
\usepackage[margin=25mm]{geometry}
\usepackage[unicode,hidelinks]{hyperref}
\setlength{\emergencystretch}{3em}
\setlength{\parskip}{0.4em}
\setlength{\parindent}{0pt}
\allowdisplaybreaks[2]
\raggedbottom
\title{\texorpdfstring{Quadratic rejection mass under a concrete phase moment}{Quadratic rejection mass under a concrete phase moment}}
\author{LangevinLean mathematical exposition}
\date{}
\begin{document}\maketitle
\noindent Source module: \texttt{Kinetic\allowbreak{}Langevin/\allowbreak{}Hamiltonian\allowbreak{}Rejection.\allowbreak{}lean}. The accompanying HTML pairs this mathematical exposition with the exact Lean statements and records the scope of the formalization.


The exact Verlet energy defect controls its actual Metropolis rejection probability. For steps between zero and one, the state-dependent bound simplifies to a single quadratic phase envelope. Integrating this bound requires precisely the velocity and force second moments displayed in the theorem, and supplies the vanishing error needed for a fixed-horizon weak limit.


\section{\texorpdfstring{Apply the scalar rejection inequality to the actual energy defect}{Apply the scalar rejection inequality to the actual energy defect}}


The acceptance is defined by the actual defect \(d=H(V_hz)-H(z)\). The already proved scalar inequality \(1-\min\{1,e^{-d}\}\le |d|\) therefore applies by substitution, with no acceptance approximation.



{\raggedright\noindent Lean declarations: \texttt{verlet\_\allowbreak{}rejection\_\allowbreak{}le\_\allowbreak{}energy\_\allowbreak{}defect}.\par}

\section{\texorpdfstring{Reduce all intermediate stages to a quadratic phase envelope}{Reduce all intermediate stages to a quadratic phase envelope}}


Put \(a=\|v\|\), \(b=\|\nabla f(x)\|\), \(w=v-h\nabla f(x)/2\), and \(q=a^2+b^2\). For \(0\le h\le1\), the triangle inequality gives \(\|w\|\le a+b\), hence \(\|w\|^2\le2q\). Also \(2\|w\|b\le\|w\|^2+b^2\), so \(\|w\|(2b+Mh\|w\|)\le(2M+3)q\). Substitution in the earlier exact defect estimate yields \[|H(V_hz)-H(z)|\le h^2\frac{M(2M+11)}8\bigl(\|v\|^2+\|\nabla f(x)\|^2\bigr).\] Lean retains the intermediate nonnegative scalar bounds, checks every multiplication by \(h\) and \(M\), and combines them by polynomial arithmetic. The following theorem composes this estimate with the scalar rejection inequality.



{\raggedright\noindent Lean declarations: \texttt{verlet\_\allowbreak{}energy\_\allowbreak{}defect\_\allowbreak{}quadratic}, \texttt{verlet\_\allowbreak{}rejection\_\allowbreak{}le\_\allowbreak{}quadratic}.\par}

\section{\texorpdfstring{Integrate the proved rejection envelope}{Integrate the proved rejection envelope}}


For any finite phase measure with integrable \(q(z)=\|v\|^2+\|\nabla f(x)\|^2\), the rejection factor is measurable and lies between zero and one, so it is integrable. The pointwise bound integrates to \[\int(1-a_h)\,d\pi\le h^2\frac{M(2M+11)}8\int q\,d\pi.\] Lean proves bounded integrability of rejection directly from the acceptance bounds and applies monotonicity of the ordinary integral against the explicitly assumed finite quadratic moment. No one-step trajectory or stationary-flow conclusion is supplied as a premise.



{\raggedright\noindent Lean declarations: \texttt{verlet\_\allowbreak{}rejection\_\allowbreak{}integral\_\allowbreak{}le}.\par}

\end{document}
